\Needspace{13\baselineskip}
\section{R9: Learned proposals with reproducible verification}\label{sec:learning}

\Q{R9.Q1}{Accuracy of carrier and helper proposals}{What accuracy do current models achieve on the exact carrier-invention task, and how should the proposed experiment be designed?}

\answer The examined sources do not provide an accuracy estimate that can be transferred to this precise Leant2 setting: a frozen Lean type, a small contract, a specified carrier/helper grammar, a fixed realization budget, and a final Lean contract proof. This is a missing measurement, not evidence of low accuracy and not a reason to abandon the approach. General synthesis or proof-premise success rates would answer a different question.

Evaluate the proposal stage and the complete system separately. A suggested carrier can parse and type-check yet be useless within the available step/finisher grammar. Conversely, a mathematically sufficient carrier can fail because its realization opens an intractable subtree. The primary endpoint should be the number of fully certified programs at equal total work or equal wall-clock budget; intermediate endpoints explain why it changes.

Use held-out task families spanning option state, tupled state, continuations, indexed motives, generalized accumulators, and invented helpers. Remove reference carriers and helper definitions from prompts. Include structural variants that change parameter order, index expressions, or examples without merely changing variable names. Keep closely related variants within one train/evaluation split to limit contamination. Report textbook tasks separately from deliberately novel combinations.

Compare at least the native grammar, model-proposed carriers/helpers, direct model-proposed programs, and an oracle-carrier condition in which the reference carrier is supplied but its realization is still synthesized. The oracle condition distinguishes a weak proposal mechanism from a weak body synthesizer. At fixed proposal counts such as $k=1,3,10$, record parsing/type validity, grammar membership, sufficient-carrier rate under a declared realization bound, checked-contract success, axiom compliance, and export/compilation success.

Include proposal latency, cold and warm index/model costs, downstream search work, and failed proposals in the totals. A high top-$k$ suggestion rate can lose overall if requesting and elaborating the suggestions consumes the interactive budget. Use paired comparisons on the same tasks and confidence intervals that respect task-family grouping. Publish ordered proposal transcripts so that downstream search can be reproduced independently of model service availability.

No model was run for this review, and no proposal-accuracy percentage is asserted. The concrete recommendation is to begin with one bounded proposal batch outside the inner loop, targeting the carrier, motive-generalization set, or helper signature. This exposes the largest possible structural benefit while keeping the experiment interpretable. Learned premise selection, represented by the current LeanHammer work, is a useful separate baseline but should not substitute for measuring this task.\cite{leanhammer}

\Q{R9.Q2}{Determinism and auditability beyond receipts}{Can the model's contribution be made deterministic enough for correctness-sensitive use, and what should be recorded?}

\answer Separate \emph{verification reproducibility}, \emph{search replay}, and \emph{model regeneration}. A closed program and contract proof can be checked without the model. A search can replay an ordered proposal transcript without regenerating it. Reproducing the same transcript from a model requires additional control over weights, tokenizer, prompt format, inference implementation, numerical behavior, batching, and tie-breaking. A fixed seed or temperature zero alone does not establish all these conditions.

For a locally controlled model, record immutable model and tokenizer identifiers, quantization and inference settings, prompt-template version, canonical serialized query, context-selection policy, and ordered outputs. For a remote service whose implementation is not fixed by the client, promise transcript replay rather than bit-for-bit regeneration. In either case, the certificate should be independent of the model's claimed confidence or explanation.

The safest proposal interface returns data: declaration names, carrier ASTs, bounded typed sketches, and explicit generalization choices. Do not execute arbitrary proposed commands, import directives, metaprograms, or tactic code in the authoritative environment. Validate the proposal grammar, node/depth limits, allowed constants, and scope before elaboration in a scratch context. Newly synthesized helper declarations should be dependency-checked and kernel-checked before becoming available to later steps.

The acceptance boundary must compare the result with the frozen query, not with a type or contract supplied by the proposer. It must also audit transitive axioms, reject escaped locals and unresolved holes, and preserve the declared execution policy. Resource limits remain necessary even for rejected proposals: a malicious or accidental oversized expression can consume excessive parsing, elaboration, or reduction resources without ever reaching the kernel gate. Optional remote prompting should have a separate, explicit policy for potentially private local definitions and contracts.

The proposal's ``only reorder, never change the set'' rule requires an actual membership test. A carrier outside the finite native grammar is not a reordering. Either translate it into a recipe already within the fixed grammar or put it in a clearly labelled extension lane with grammar $\G^+$. A native fair fallback preserves eventual exploration under its stated assumptions, not identical finite-deadline results. An extension lane cannot justify an impossibility claim about the original or expanded language merely because it found nothing.

For auditability, a receipt should contain the original query and environment digest, search-grammar and rank-model versions, proposal transcript, exact universe treatment, work ledger, closed program, exact contract proof, transitive axiom inventory, and export relationship. Record unknowns, resource cuts, and the origin of each proposed alternative. A replay command should verify the certificate before optionally replaying search; correctness verification must not require repeating a statistical choice.

The default should remain Lean-native and model-free. An optional model can be an effective source of hypotheses, especially for carriers and helpers, but the same proposal interface should also serve deterministic grammars and user sketches. This makes model guidance replaceable and keeps the research question about useful proposals separate from the trust question about accepted results.
